\documentclass{article}
\usepackage[UTF8]{ctex}
\usepackage{amsmath,amssymb,amsfonts}
\usepackage{enumitem}
\usepackage{geometry}
\geometry{a4paper, margin=1in}
% === 错题本排版优化（不参与题目标识） ===
\setlength{\parindent}{0pt}
\setlength{\parskip}{0.35em plus 0.10em minus 0.08em}
\setlength{\emergencystretch}{3em}
\setlist{itemsep=0.25em, topsep=0.40em, parsep=0pt}
\allowdisplaybreaks[2]
% === 错题本排版优化结束 ===
\begin{document}
\section*{上海财经大学2026年高等代数真题}
\section*{45．上海财经大学2026年高等代数试卷}
一．（10 分）令 \(f_{1}(x)=x^{2}+3 x+2, f_{2}(x)=x^{4}-2 x^{3}-3 x^{2}+4 x+4\) ，求它们的极大公因式 \(d(x)\) ，并求多项式 \(u_{1}(x), u_{2}(x)\) ，使得 \(f_{1}(x) u_{1}(x)+f_{2}(x) u_{2}(x)=d(x)\) ．

二．（15 分）计算下列矩阵的行列式：

\[
A=\left(\begin{array}{ccccc}
2 & x+1 & x+1 & \cdots & x+1 \\
2 & x+2 & x & \cdots & x \\
2 & x & x+3 & \cdots & x \\
\vdots & \vdots & \vdots & & \vdots \\
2 & x & x & \cdots & x+n
\end{array}\right) .
\]

三．（15 分）设 \(a, b\) 是常数，且有下列线性方程组

\[
\left\{\begin{array}{l}
x_{1}+3 x_{2}-a x_{3}=2, \\
x_{1}+4 x_{2}+a x_{3}=3+2 b, \\
3 x_{1}+11 x_{2}+(3-2 a) x_{3}=8+b .
\end{array}\right.
\]

1．\(a, b\) 为何值时，方程组有解？\\

2．当方程组有解时，求出所有解。

四．（10 分）设 \(A, B\) 是 \(n\) 阶方阵，满足 \(A B=B A\) ，且 \(A B\) 满足多项式 \(f(x)=x^{2025}+x+1\) ，求证： \(I+A B\) 可逆．

五．（15 分）设 \(M_{n}(\mathbb{R})\) 是实数域 \(\mathbb{R}\) 上 \(n\) 阶矩阵全体。

\begin{itemize}

  \item[1．] 令 \(W\) 是 \(M_{n}(\mathbb{R})\) 中的第一行元素全为 0 的矩阵构成的集合；求证：\(W\) 是 \(M_{n}(\mathbb{R})\) 的子空间。

  \item[2．] 设 \(V\) 是 \(M_{n}(\mathbb{R})\) 的子空间，且 \(V\) 中的非零矩阵可逆，求证：\(V\) 的维数小于等于 \(n\) ．
\end{itemize}

六．（15 分）假设 \(a, b, c\) 是不全为零的实数，已知二次型 \(f\left(x_{1}, x_{2}, x_{3}\right)=\) \(\left(x_{1}+x_{2}+x_{3}\right)\left(a x_{1}+b x_{2}+c x_{3}\right)\) ．

\begin{itemize}

  \item[1．] 求证：\(f\left(x_{1}, x_{2}, x_{3}\right)\) 的秩等于 2 当且仅当 \(a, b, c\) 不全相等。

  \item[2．] 当 \(f\left(x_{1}, x_{2}, x_{3}\right)\) 的秩等于 2 时，求 \(f\left(x_{1}, x_{2}, x_{3}\right)\) 的正、负惯性指数。
\end{itemize}

七．（15 分）设 \(A\) 是 \(n\) 阶实方阵，\(A\) 和 \(A^{*}\) 都满足多项式 \(f(x)=x^{3}-x^{2}-x+1\) ，且 \(|A|<0\) 。\\

1 ．求 \(|A|\) 。\\

2．求证：矩阵 \(A\) 可以对角化。

八．（15 分）设 \(V=M_{n}(\mathbb{R})\) 是实数域 \(\mathbb{R}\) 上 \(n\) 阶矩阵全体，定义 \(\sigma: V \times V \rightarrow \mathbb{R}, \sigma(A, B)=\operatorname{tr}\left(A B^{\mathrm{T}}\right)\) 。\\

1．证明：\(V\) 是 \(\sigma\) 下构成的一个 \(\mathbb{R}\) 上的内积空间。\\

2．设 \(V_{1}\) 是由全体对称矩阵构成的子空间，\(V_{2}\) 是由全体反对称矩阵构成的子空间，证明：\(V_{1}\) 是 \(V_{2}\) 的正交补空间。

九．（20 分）设 \(A=\left(\begin{array}{cccc}1 & -1 & 0 & 1 \\ 1 & 1 & 1 & 0 \\ 0 & -1 & 1 & 1 \\ 1 & 0 & 1 & 1\end{array}\right)\)\\

1．求 \(A\) 的 Jordan 标准型 J。\\

2 ．求可逆矩阵 \(P\) ，使得 \(P^{-1} A P=J\) 。\\

3．求 \(A\) 的极小多项式。

十．（20 分）设 \(\varphi\) 是数域 \(\mathbb{F}\) 上 \(n\) 维线性空间 \(V\) 上的线性变换，其特征多项式为 \(f(\lambda)\) ．证明以下 3 个结论等价：\\
（1）\(V\) 只有平凡的 \(\varphi\)－不变子空间．\\
（2）对任意非零向量 \(\alpha \in V\) ，都有 \(V=L\left(\alpha, \varphi(\alpha), \varphi^{2}(\alpha), \cdots\right)\) ．\\
（3）\(f(\lambda)\) 是数域 \(\mathbb{F}\) 上的不可约多项式．
\end{document}
